% CLASS 5 -- the last subsection of the notes' section 1.4, sec.idealisedAgainstHawkes.
% Taught on Monday 19 October, after section 1.5. Equation eq.idealisedUpdate is exact for the
% book of the CKS placement. The class drives that book and the Hawkes flow, placed
% geometrically, with the same events from one opening book, and asks how far the identity
% carries over: first through the time at which the two mid-prices part, then through their
% discrepancy after it, with orders of one share and of random size. 30 frames.
%
% SWITCH POINTS to notebook 10-the-idealised-book-against-the-hawkes-book, each signposted in
% the \note of the frame before it:
%   after "The standing assumption"                -> 1. Three flows on one path
%   after "Beyond the corollary, L orders ..."     -> 2. When the mid-prices part
%   after "The variance of the discrepancy, ..."   -> 3. The discrepancy: mean and variance
%   after "What random sizes change"               -> 4. Uncertain sizes
%
% Overrunning, cut in this order: the pictures of the drawdown and of the touches on its path
% (said over the touch walks), the mirror (said over the moments), then the table of the two
% placements (said over the CKS placement).
%
% THE FRAME CONTRACT. These slides are projection material, not reference material: the
% students read the notes, and they only ever see a frame while it is being talked over.
% So every frame carries ONE OBJECT -- a display, a statement of at most two lines, a TikZ
% picture, a table or an algorithmic block -- and under it at most three CAPTIONS. A caption
% is a gloss (`term: noun phrase'), a fragment of subject and predicate, or a condition
% written as mathematics. It is never a sentence, never a subordinate clause, and never a
% repetition of the frame title: what the title says, the body does not say again, and what
% follows from the object is spoken. At most 12 words of prose each, 25 to a frame, hard
% cap 40. A statement states its head claim and stops; its consequence is spoken.
% What is said over the object goes in \note{}, which \documentclass[script]{unito26slides}
% prints beside each frame and the deck itself never shows. Check with
%   python3 .claude/skills/writing-in-tex/scripts/frame_density.py <file>
% Open the class on the chapter roadmap, as classes 1 and 2 do. Do not close on a summary
% frame: a slide that recites what the class just did is the spoken half written down.
%
% The CKS placement is shown as a bare `algorithmic' environment, as class 2 shows the fold:
% `algorithm' is a float, and floats do not place inside a frame. A lemma of the notes is
% beamer's own `lemma', on beamer's counter.
% The event G of the notes, on which every statement is made, is stated once, on the frame of
% the standing assumption, and no statement carries it.
%
% LABELS. A label is defined once across the class files, and a cross-file \ref prints ??
% under \includeonly. Class 5 owns: def.threeBooks, assumption.balancedFlow, def.partingTime,
% lemma.oneBestPrice, corol.partingAtOneQuote, prop.partingProbability, corol.quietBook,
% prop.cksTouchWalks, prop.hawkesTouch, prop.fluidTouch, corol.twoRegimes, def.discrepancy,
% prop.discrepancyMoments.
% remark.wholeBookImbalance survives in the \note of the standing assumption, without its label.
% The paragraphs of the notes that precede prop.fluidTouch and follow corol.twoRegimes and
% prop.discrepancyMoments are unnumbered, and so are their four frames.
% assumption.idealisedBook, prop.idealisedUpdate, def.orderFlowImbalance,
% def.orderFlowProcess, prop.signedSizeContribution, def.directionSymmetry and
% prop.forecastAtLongHorizons belong to class 4; example.idealisedPlacement and
% def.geometricPlacement carry no label anywhere in the deck. Each is RESTATED HERE WITHOUT
% LABELS.
%
% A frame holds 218pt of body -- eighteen 12pt lines, or seven two-line paragraphs.

\subsection{The idealised flow against the Hawkes flow}

\begin{frame}
  \frametitle{Outline}
% One \nocite per line -- a multi-key \nocite spanning lines breaks bibtex.
\nocite{CKS14pri}
\nocite{CST10sto}
\nocite{BDHM13som}
\begin{center}
\begin{tikzpicture}[font=\scriptsize, node distance=3mm,
    stage/.style={draw=gray!60, rounded corners=2pt, align=center, inner sep=4pt,
                  text width=1.45cm, fill=gray!8},
    past/.style={stage, fill=gray!20, draw=gray!70},
    current/.style={stage, fill=blue!15, draw=orange!70, line width=1pt}]
  \node[past] (a) {1.1\\the book};
  \node[past, right=of a] (b) {1.2\\the fold};
  \node[past, right=of b] (c) {1.3\\point processes};
  \node[current, right=of c] (d) {1.4\\price formation};
  \node[past, right=of d] (e) {1.5\\recorded data};
  \draw[-{Stealth}, gray!60] (a) -- (b);
  \draw[-{Stealth}, gray!60] (b) -- (c);
  \draw[-{Stealth}, gray!60] (c) -- (d);
  \draw[-{Stealth}, gray!60] (d) -- (e);
\end{tikzpicture}
\end{center}

The first time the mid-prices part

Their discrepancy after it

Orders of random size
\note{
We return to section 1.4 after the recorded data of section 1.5. Section 1.4 wrote the change
of the mid-price as the order flow imbalance over $2\uniformDepth$, plus a residual of at most a
tick, under an assumption on the book. The book that the Hawkes flow of the simulator builds
does not satisfy it. This class measures, on that flow, how far the identity carries over: first
the time at which the mid-prices of the two books first part, then their discrepancy after it,
then what changes when the orders carry random sizes.
}
\end{frame}

\begin{frame}
  \frametitle{For the CKS book the update holds on every path}
\[
 \Delta\midPrice_{t,\ofiWindow}
 = \frac{\tickSizeOfLOB}{2\uniformDepth}\OFI_{t,\ofiWindow} + \residual_{t,\ofiWindow},
 \qquad
 \vert\residual_{t,\ofiWindow}\vert \leq \tickSizeOfLOB
\]

$\cks$: our version of the book of Cont, Kukanov and Stoikov

$\OFI^{\cks}_{t,t} = \pressure^{\transpose}\multiCountingProc(t)$: the signed count

The Hawkes flow, placed geometrically: the same events, one opening book
\note{
The update of class 4, under the idealised book. The placement of the next frame satisfies the
assumption on every path, so for its book, which we call the CKS book, the update holds event by
event, and every event contributes its pressure: the imbalance is the signed count.
The simulator places the same events geometrically, and its book satisfies the assumption at no
value of $\offsetParameter$. We fold both from one opening book with the same events and ask how
far the identity carries over to the second.
}
\end{frame}

\begin{frame}
  \frametitle{From a Hawkes event to a CKS order}
{\footnotesize
\begin{algorithmic}
 \REQUIRE the book before $\arrivalTimes[]_n$, and the type $\event[n]$, of direction $d$
 \IF{$\event[n]$ is a market order}
 \STATE submit one share at the touch of the side $-d$
 \ELSIF{$\event[n]$ is a withdrawal}
 \STATE cancel one share at the touch of the side $d$
 \ELSIF{the touch of the side $d$ holds less than $\uniformDepth$}
 \STATE submit one share at that touch
 \ELSE
 \STATE submit one share at that touch improved by one tick
 \ENDIF
\end{algorithmic}
}

$\LOBspread = \tickSizeOfLOB$: the improved price is the opposite touch, a trade

$\LOBspread > \tickSizeOfLOB$: the improved price is a free tick, a new touch

Nothing drawn beyond the event: $\cks$ is $\sigmaAlgebra^{\multiCountingProc}$-measurable

The placement of class 4 under which the idealised book holds
\note{
Each event of the Hawkes process becomes one message of one share, read against the book as it
stands, and the matching is the ordinary matching of the book. A market order is a limit order of
one share at the opposite touch, so its market part is its whole size. A limit order at a full
touch is priced one tick better than that touch, and the matching decides what it does: at a
spread of one tick that price is the opposite touch, where it trades one share, as a market order
would; at a wider spread the tick is free, and the order opens a new touch there.
The sizes are units because an order of two shares may reach past a touch that holds one, against
part ii of the assumption. $\sigmaAlgebra^{\multiCountingProc}$ is the $\sigma$-algebra of the
arrival times and the types: the CKS book is a function of the opening book and of the path.
}
\end{frame}

\begin{frame}
  \frametitle{The CKS and the geometric placement of one event}
\begin{center}
{\footnotesize
\begin{tabular}{@{}lp{3.3cm}p{4.0cm}@{}}
\toprule
event & CKS & geometric \\
\midrule
market buy & one share at $\bestAskPrice$ & $\orderSize_n$ shares, walking the ask \\
limit buy & one share at $\bestBidPrice$, or one tick above it once full
 & $\orderSize_n$ shares at $\bestAskPrice - \tickSizeOfLOB(1 + \priceOffset_n)$ \\
withdrawal, bid & one share from $\bestBidPrice$
 & $\orderSize_n \wedge \size^{b}(p_n)$ from $p_n$, drawn by size \\
\bottomrule
\end{tabular}
}
\end{center}

$\Prob(\priceOffset_n = k) = \offsetParameter(1-\offsetParameter)^{k}$, independent of the flow

$\orderSize_n$: one share, or lognormal of mean $\meanSize$
\note{
The sell types are the mirror image. The first column is the algorithm of the last frame. The
second is the geometric placement of class 4: a market order walks the opposite side; a limit
order is priced $1 + \priceOffset_n$ ticks from the opposite best quote, so it rests in full on
its own side, behind its touch when $\priceOffset_n \geq \LOBspread/\tickSizeOfLOB$; a
withdrawal draws an occupied price of its side in proportion to the size resting there, and
cancels its size or what rests there, whichever is less.
Under the CKS placement a limit order may initiate a transaction. Under the geometric placement
it never does.
}
\end{frame}

\begin{frame}
  \frametitle{Three books on one path}
\begin{defi}
\label{def.threeBooks}
 From one opening book of spread one tick, $L$ levels a side, touches
 $\nthBestBidSize[1]_0, \nthBestAskSize[1]_0 \leq \uniformDepth$ and $\uniformDepth$ behind,
 fold:\par
 \hfill\begin{tabular}{lll}
  $\cks$ & CKS & $(\arrivalTimes[]_n, \event[n])$ \\
  $\hawOne$ & geometric & $(\arrivalTimes[]_n, \event[n], 1)$ \\
  $\hawQ$ & geometric & $(\arrivalTimes[]_n, \event[n], \orderSize_n)$,
  opening sizes $\times\,\meanSize$
 \end{tabular}\hfill\null
\end{defi}

$(\arrivalTimes[]_n, \event[n])_n$: the arrival times and types of the Hawkes process

$\orderSize_n$: the lognormal sizes of the order flow, of mean $\meanSize$
\note{
The opening book has one level on every tick, its touches hold at most $\uniformDepth$ and every
level behind them $\uniformDepth$, so it satisfies part iii of the assumption. $\hawOne$ is the Hawkes flow with orders of one share and the geometric rule
for prices; $\hawQ$ gives the orders the sizes of the order flow of class 4, lognormal of mean
$\meanSize$, which act as an uncertainty in the size of an observed order. A quantity of a book carries its superscript, as in
$\midPriceOf{\cks}_t$ or $\nthBestBidSize[1,\hawOne]_t$, and a quantity indexed by
$\arrivalTimes[]_n$ is read after the $n$-th event.
The geometric placement draws for each event an offset or a level, and $\hawQ$ also a size. Given
$\sigmaAlgebra^{\multiCountingProc}$ and the draws of the earlier events, these have the laws of
the geometric placement, the level of a withdrawal being drawn on the book as the earlier events
left it. We represent that draw by $U_n$, uniform on $(0,1)$ and independent of the path and of
the other draws: the withdrawal takes the first level of its side, counted from the touch
outwards, at which the cumulated size exceeds $U_n$ times the size of the side.
}
\end{frame}

\begin{frame}
  \frametitle{Each step changes one thing}
\begin{center}
\begin{tikzpicture}[font=\scriptsize, node distance=1.5cm,
    book/.style={draw=gray!60, rounded corners=2pt, align=center, inner sep=4pt,
                 text width=1.9cm, fill=gray!8}]
  \node[book, fill=blue!12] (c) {$\cks$\\[2pt] CKS placement\\one share};
  \node[book, right=of c] (h) {$\hawOne$\\[2pt] geometric placement\\one share};
  \node[book, right=of h, fill=orange!15] (q)
       {$\hawQ$\\[2pt] geometric placement\\$\orderSize_n$ shares};
  \draw[-{Stealth}, gray!70] (c) -- node[above, align=center] {where the\\orders go} (h);
  \draw[-{Stealth}, gray!70] (h) -- node[above, align=center] {the size of\\each order} (q);
\end{tikzpicture}
\end{center}

$\cks$ against $\hawOne$: the cost of placing orders away from the touch

Against $\hawQ$: also the cost of not seeing each size
\note{
The CKS book and $\hawOne$ receive the same events and the same sizes, and differ only in where
the orders go. $\hawOne$ and $\hawQ$ receive the same events and use the same placement, and
differ only in the size of each order. The update of the CKS book counts events, by the
signed-size contribution of class 4. So the comparison with $\hawOne$ measures the cost of
placing the orders away from the touch, and the comparison with $\hawQ$ adds the cost of not
seeing the size of each order.
}
\end{frame}

\begin{frame}
  \frametitle{A limit buy at a spread of $k \geq 2$ ticks}
\begin{center}
\begin{tikzpicture}[x=0.34cm, y=0.24cm, font=\scriptsize,
    bid/.style={fill=blue!40, draw=blue!60}, ask/.style={fill=orange!40, draw=orange!65},
    span/.style={{Stealth[length=3pt]}-{Stealth[length=3pt]}, gray!50!black},
    moves/.style={blue!65, thick}, stays/.style={gray!55, thick}]
  \begin{scope}
    \fill[gray!10] (5.5,0) rectangle (8.5,12);
    \booksizeaxis{0.4}{12.6}{6}{2}
    \draw[gray!55] (0.4,0) -- (12.6,0);
    \draw[gray!45] (0.4,7.5) -- (12.6,7.5);
    \foreach \x in {1,...,5} \draw[bid] (\x-0.3,0) rectangle (\x+0.3,4);
    \draw[ask] (9-0.3,0) rectangle (9+0.3,2);
    \foreach \x in {10,11,12} \draw[ask] (\x-0.3,0) rectangle (\x+0.3,4);
    \node[above, inner sep=1pt] at (5,4) {$\uniformDepth$};
    \draw[moves] (6,7.5) -- (6,11.5);
    \fill[blue!65] (6,11.5) circle (1.5pt) node[right] {$1$};
    \node[below] at (5,-0.2) {$\bestBidPrice$};
    \node[below] at (9,-0.2) {$\bestAskPrice$};
    \draw[span] (5,-3.6) -- node[below] {$k$ ticks} (9,-3.6);
    \node at (6.5,13.6) {$\cks$};
  \end{scope}
  \begin{scope}[xshift=5.3cm]
    \fill[gray!10] (5.5,0) rectangle (8.5,12);
    \booksizeaxis{0.4}{12.6}{6}{2}
    \draw[gray!55] (0.4,0) -- (12.6,0);
    \draw[gray!45] (0.4,7.5) -- (12.6,7.5);
    \foreach \x/\h in {1/3, 2/5, 3/2, 4/4, 5/2} \draw[bid] (\x-0.3,0) rectangle (\x+0.3,\h);
    \foreach \x/\h in {9/2, 10/4, 11/1, 12/3} \draw[ask] (\x-0.3,0) rectangle (\x+0.3,\h);
    \foreach \x/\h in {8/1.6, 7/0.96, 6/0.576} {
      \draw[moves] (\x,7.5) -- (\x,7.5+\h);
      \fill[blue!65] (\x,7.5+\h) circle (1.5pt);}
    \foreach \x/\h in {5/0.346, 4/0.207, 3/0.124, 2/0.075} {
      \draw[stays] (\x,7.5) -- (\x,7.5+\h);
      \fill[gray!55] (\x,7.5+\h) circle (1.5pt);}
    \node[right] at (8,9.1) {$\offsetParameter$};
    \node[below] at (5,-0.2) {$\bestBidPrice$};
    \node[below] at (9,-0.2) {$\bestAskPrice$};
    \draw[span] (5,-3.6) -- node[below] {$k$ ticks} (9,-3.6);
    \node at (6.5,13.6) {$\hawOne$};
  \end{scope}
\end{tikzpicture}
\end{center}

Stems: the law of the price of the incoming limit buy

$\cks$: one tick inside if its touch holds $\uniformDepth$, else at it

$\hawOne$: inside iff $\tickSizeOfLOB(1+\priceOffset_n) < \LOBspread$, with probability
$1-(1-\offsetParameter)^{k-1}$
\note{
A limit buy, the sell being the mirror, at a spread of $k = 4$ ticks, the free ticks shaded.
The CKS book sends it to its touch while the touch holds less than $\uniformDepth$, and one tick
inside once the touch is full, where it opens a new touch of one share.
The geometric placement prices it $1 + \priceOffset_n$ ticks below the best ask, whatever the
queue. It improves the bid by $j = k - 1 - \priceOffset_n$ ticks, $1 \leq j \leq k - 1$, with
probability $\offsetParameter(1-\offsetParameter)^{k-1-j}$, joins the touch with probability
$\offsetParameter(1-\offsetParameter)^{k-1}$, and rests behind it with probability
$(1-\offsetParameter)^{k}$. At a full CKS touch the two books improve the bid by the same tick
with probability $\offsetParameter(1-\offsetParameter)^{k-2}$: the rows $k \geq 2$ of the parting
table to come.
}
\end{frame}

\begin{frame}
  \frametitle{A limit buy at a spread of one tick}
\begin{center}
\begin{tikzpicture}[x=0.34cm, y=0.24cm, font=\scriptsize,
    bid/.style={fill=blue!40, draw=blue!60}, ask/.style={fill=orange!40, draw=orange!65},
    stays/.style={gray!55, thick}]
  \begin{scope}
    \booksizeaxis{0.4}{10.6}{6}{2}
    \draw[gray!55] (0.4,0) -- (10.6,0);
    \draw[gray!45] (0.4,7.5) -- (10.6,7.5);
    \foreach \x in {1,...,5} \draw[bid] (\x-0.3,0) rectangle (\x+0.3,4);
    \draw[ask] (6-0.3,0) rectangle (6+0.3,1);
    \draw[orange!65, dashed] (6-0.3,1) rectangle (6+0.3,2);
    \foreach \x in {7,...,10} \draw[ask] (\x-0.3,0) rectangle (\x+0.3,4);
    \node[above, inner sep=1pt] at (5,4) {$\uniformDepth$};
    \draw[orange!80!black, thick] (6,7.5) -- (6,11.5);
    \fill[orange!80!black] (6,11.5) circle (1.5pt) node[right] {$1$};
    \draw[-{Stealth}, orange!80!black] (6,7.3) -- (6,2.3);
    \node[right, inner sep=1pt] at (6.3,5.4) {trades};
    \node[below left, inner sep=1pt] at (5.3,-0.2) {$\bestBidPrice$};
    \node[below right, inner sep=1pt] at (5.7,-0.2) {$\bestAskPrice$};
    \node at (5.5,13.6) {$\cks$};
  \end{scope}
  \begin{scope}[xshift=4.9cm]
    \booksizeaxis{0.4}{10.6}{6}{2}
    \draw[gray!55] (0.4,0) -- (10.6,0);
    \draw[gray!45] (0.4,7.5) -- (10.6,7.5);
    \foreach \x/\h in {1/3, 2/5, 3/2, 4/4, 5/2} \draw[bid] (\x-0.3,0) rectangle (\x+0.3,\h);
    \foreach \x/\h in {6/2, 7/4, 8/1, 9/3, 10/2} \draw[ask] (\x-0.3,0) rectangle (\x+0.3,\h);
    \foreach \x/\h in {5/1.6, 4/0.96, 3/0.576, 2/0.346, 1/0.207} {
      \draw[stays] (\x,7.5) -- (\x,7.5+\h);
      \fill[gray!55] (\x,7.5+\h) circle (1.5pt);}
    \node[right] at (5,9.1) {$\offsetParameter$};
    \node[below left, inner sep=1pt] at (5.3,-0.2) {$\bestBidPrice$};
    \node[below right, inner sep=1pt] at (5.7,-0.2) {$\bestAskPrice$};
    \node at (5.5,13.6) {$\hawOne$};
  \end{scope}
\end{tikzpicture}
\end{center}

$\cks$, touch at $\uniformDepth$: one share traded at $\bestAskPrice$

$\hawOne$: at the touch with probability $\offsetParameter$, behind it otherwise

$\cks$ moves the best ask iff that touch held one share
\note{
At a spread of one tick no tick is free. The CKS limit buy at a full bid is priced one tick above
it, at the best ask, and trades one share there, as a market buy would: the bid stays full, the
ask touch loses a share, and the best ask moves up a tick when that share was its last. Below
$\uniformDepth$ the order joins the bid. The geometric limit buy is priced $\priceOffset_n$ ticks
below the best bid, at the touch when $\priceOffset_n = 0$ and behind it otherwise, and it moves
no best price. So at one tick no limit order improves a touch, and a best price moves only when
its touch loses its last share: the quiet book, to come.
}
\end{frame}

\begin{frame}
  \frametitle{The standing assumption}
\begin{assumption}
\label{assumption.balancedFlow}
 Direction-symmetric, $\branchingRatio < 1$, balanced per side:
 $\stationaryIntensity[3] = \stationaryIntensity[2] + \stationaryIntensity[5]$ and
 $\stationaryIntensity[4] = \stationaryIntensity[1] + \stationaryIntensity[6]$.
 The opening best bid lies above one tick.
\end{assumption}

$\swapMatrix\stationaryIntensity = \stationaryIntensity$: the two conditions are one

$\size^{b,\hawOne}_t(\tickSizeOfLOB\N) = \size^{b}_0(\tickSizeOfLOB\N)
+ (\sideVector{b})^{\transpose}\multiCountingProc(t)$: a walk of bounded mean

Statements on $G$: no empty side or floor, on path and mirror
\note{
Throughout, we work on the event $G$ that, on the horizon considered, on the path and on its
mirror, no side of any of the three books is empty and no limit price meets the floor of the
geometric placement. Every statement of the class is made on $G$, and expectations and laws are
taken on it. $G$ changes the laws of the draws only by the probability of its complement, small
on a book of many levels priced far above one tick, and it is invariant under the exchange of the
path with its mirror, the path $(\arrivalTimes[]_n, \swapMatrix(\event[n]))_n$ with the same draws.
With orders of one share a withdrawal removes one share wherever it is drawn, so the whole bid of
$\hawOne$ is the opening bid, $(L-1)\uniformDepth$ plus its touch, plus the limit buys less the
market sells and the withdrawals from the bid: $\sideVector{b} = \unitVector_3 - \unitVector_2 - \unitVector_5$, likewise
$\sideVector{a}$, and $\sideVector{b} - \sideVector{a} = \pressure$. Under balance the mean of the
walk is bounded and its variance grows as
$t\,(\sideVector{b})^{\transpose}\countCovarianceRate\sideVector{b}$, so a side lasts a time of
the order of $(L\uniformDepth)^2/(\sideVector{b})^{\transpose}\countCovarianceRate\sideVector{b}$,
the life of a side.
The difference of the two sides, the whole-book imbalance of $\hawOne$, is that of the opening
touches plus the signed count, whatever the draws. The CKS book has the same, and turns it into a mid-price at the fixed depth
$\uniformDepth$ per tick; $\hawOne$ has no fixed depth.
Now to notebook 10, section 1, Three flows on one path.
}
\end{frame}

\begin{frame}
  \frametitle{The mid-prices part at $\partingTime$}
\begin{defi}
\label{def.partingTime}
 $\partingTime := \inf\{t \geq 0 : \midPriceOf{\cks}_t \neq \midPriceOf{\hawOne}_t\}$.
\end{defi}

Before $\partingTime$: the same mid-prices, possibly different queues

For $\hawQ$: the same, with $\hawQ$ in place of $\hawOne$
\note{
The parting time compares the mid-prices only. The queues of the two books may differ long
before it: a limit order that the geometric placement puts behind the touch, or a withdrawal
drawn behind it, changes a queue of $\hawOne$ and no best price.
}
\end{frame}

\begin{frame}
  \frametitle{Each event moves one best price at most}
\begin{lemma}
\label{lemma.oneBestPrice}
 In each book an event moves at most one best price, in the direction
 $\pressure_{\event[n]}$.
 At agreeing best prices, $\cks$ and $\hawOne$ never move different ones.
\end{lemma}

$\pressure_{\event[n]} = +1$: up; $\pressure_{\event[n]} = -1$: down

CKS limit order, full touch, spread one tick: the opposite price
\note{
Take a buy type; the sell types are the same mutatis mutandis. A market buy trades at the best
ask and changes nothing else: one share in the CKS book and in $\hawOne$, $\orderSize_n$
shares walking the ask in $\hawQ$. Consuming the lowest asks can only raise the best ask.
A geometric limit buy is priced below the best ask and rests in full on the bid: adding size can
only raise the best bid or leave it. The CKS limit buy joins the best bid, or opens a level one
tick above it, or, at a full best bid facing a spread of one tick, trades one share at the best
ask; in the last case only the best ask can move, and up. A withdrawal from the bid removes size
from the bid alone, and can only lower the best bid.
For the second claim, the only event that moves different prices in the two books would be that
last CKS limit order. At a spread of one tick the geometric limit buy is priced at or behind the
best bid, and it moves no best price.
}
\end{frame}

\begin{frame}
  \frametitle{The mid-prices part at one quote}
\begin{corol}
\label{corol.partingAtOneQuote}
 $\partingTime$ is the first time a best price of $\cks$ differs from that of $\hawOne$,
 and then exactly one of the two best prices differs.
\end{corol}

From the common opening book, event by event
\note{
The opening book is common. If the best prices agree before an event, by the lemma at most one
of them differs after it, and if one does, the mid-prices differ. If they agree, so do the
mid-prices. The books therefore part one best bid apart, or one best ask apart, and never both.
}
\end{frame}

\begin{frame}
  \frametitle{What parts the mid-prices --- 1 of 2}
\begin{center}
\begin{tikzpicture}[x=0.34cm, y=0.24cm, font=\scriptsize,
    bid/.style={fill=blue!40, draw=blue!60}, ask/.style={fill=orange!40, draw=orange!65},
    span/.style={{Stealth[length=3pt]}-{Stealth[length=3pt]}, gray!50!black}]
  \begin{scope}
    \booksizeaxis{0.4}{11.6}{6}{2}
    \draw[gray!55] (0.4,0) -- (11.6,0);
    \foreach \x in {1,...,5} \draw[bid] (\x-0.3,0) rectangle (\x+0.3,4);
    \draw[bid] (6-0.3,0) rectangle (6+0.3,2);
    \draw[ask] (8-0.3,0) rectangle (8+0.3,3);
    \foreach \x in {9,10,11} \draw[ask] (\x-0.3,0) rectangle (\x+0.3,4);
    \draw[gray!60] (6,2) -- (6,7.2) node[above, inner sep=1pt, text=black] {$\size^{\cks}$};
    \draw[span] (6,-3.6) -- node[below] {$k$ ticks} (8,-3.6);
    \node[below] at (6,-0.2) {$\bestBidPrice$};
    \node[below] at (8,-0.2) {$\bestAskPrice$};
    \node at (6,11) {$\cks$};
  \end{scope}
  \begin{scope}[xshift=4.6cm]
    \booksizeaxis{0.4}{11.6}{6}{2}
    \draw[gray!55] (0.4,0) -- (11.6,0);
    \foreach \x/\h in {1/5, 2/2, 3/6, 4/3} \draw[bid] (\x-0.3,0) rectangle (\x+0.3,\h);
    \draw[bid] (6-0.3,0) rectangle (6+0.3,1);
    \foreach \x/\h in {8/2, 9/5, 10/1, 11/4} \draw[ask] (\x-0.3,0) rectangle (\x+0.3,\h);
    \node[above] at (6,1) {$\size^{\hawOne}$};
    \node[below] at (6,-0.2) {$\bestBidPrice$};
    \node[below] at (8,-0.2) {$\bestAskPrice$};
    \draw[span] (4,-3.6) -- node[below] {$g$} (6,-3.6);
    \draw[span] (0.7,-7.4) -- node[below] {$\size^{\hawOne}(\tickSizeOfLOB\N)$} (6.3,-7.4);
    \node at (6,11) {$\hawOne$};
  \end{scope}
\end{tikzpicture}
\end{center}

Side of the event: the side its CKS message acts on

$g$: ticks from the touch of $\hawOne$ to the next occupied price
\note{
The side of the event is the side on which its CKS message acts: the opposite side for a market
order, and for a limit order at a full touch facing a spread of one tick; its own side otherwise.
The bid is drawn as the side of the event, as for a market sell or a withdrawal from the bid. On
it $\size^{\cks}$ and $\size^{\hawOne}$ are the sizes at the two touches,
$\size^{\hawOne}(\tickSizeOfLOB\N)$ is the size resting on the whole side of $\hawOne$, and $g$ the
number of ticks from the touch of $\hawOne$ to the next occupied price behind it.
Behind its touch the CKS book holds $\uniformDepth$ on every tick, so it moves a best price by one
tick or not at all. $\hawOne$ may move it by $g$ ticks, and a limit order that improves its touch
by any number of ticks inside the spread.
}
\end{frame}

\begin{frame}
  \frametitle{What parts the mid-prices --- 2 of 2}
\begin{prop}
\label{prop.partingProbability}
 At agreeing best prices, $\LOBspread = k\tickSizeOfLOB$, given
 $\sigmaAlgebra^{\multiCountingProc}$ and the earlier draws:\par\smallskip
 {\footnotesize\hfill\begin{tabular}{@{}llll@{}}
  \toprule
  event & & $j^{\cks}$ & parting \\
  \midrule
  market & $\size^{\cks} \geq 2$ & $0$ & $\indicator{\size^{\hawOne} = 1}$ \\
   & $\size^{\cks} = 1$ & $1$ & $1 - \indicator{\size^{\hawOne} = 1, \, g = 1}$ \\
  limit, $k = 1$ & own touch $< \uniformDepth$ & $0$ & $0$ \\
   & own touch $= \uniformDepth$ & $\indicator{\size^{\cks} = 1}$
   & $\indicator{\size^{\cks} = 1}$ \\
  limit, $k \geq 2$ & $\size^{\cks} < \uniformDepth$ & $0$
   & $1 - (1-\offsetParameter)^{k-1}$ \\
   & $\size^{\cks} = \uniformDepth$ & $1$
   & $1 - \offsetParameter(1-\offsetParameter)^{k-2}$ \\
  withdrawal & $\size^{\cks} \geq 2$ & $0$
   & $\indicator{\size^{\hawOne} = 1}/\size^{\hawOne}(\tickSizeOfLOB\N)$ \\
   & $\size^{\cks} = 1$ & $1$
   & $1 - \indicator{\size^{\hawOne} = 1, \, g = 1}/\size^{\hawOne}(\tickSizeOfLOB\N)$ \\
  \bottomrule
 \end{tabular}\hfill\null\par}
\end{prop}

$j^{\cks}$, $j^{\hawOne}$: the ticks each book moves the side of the event

Parting: the probability that event $n$ parts the mid-prices

Market order: a withdrawal that draws the touch
\note{
Each book moves the best price of the side of the event by $j^{\cks}$ and $j^{\hawOne}$ ticks,
and no other best price, and the event parts the mid-prices exactly when the two differ. For a
limit order the own touch is the touch of the CKS book on the side of the order.
The CKS book holds $\uniformDepth$ on every tick behind its touch, so $j^{\cks}$ is zero or one. A
market order or a withdrawal moves the touch it takes from exactly when that touch holds one
share. A limit order at $k \geq 2$ opens a level one tick inside its touch exactly when that touch
is full. At $k = 1$ it joins its touch if that is not full, and otherwise trades one share at the
opposite touch, the touch of the side of the event, which moves exactly when it holds one share.
In $\hawOne$ a withdrawal moves the touch exactly when it draws it, with probability
$\size^{\hawOne}/\size^{\hawOne}(\tickSizeOfLOB\N)$, and the touch held one share; then by $g$.
In both books a market order is a withdrawal that draws the touch for sure: its rows are those of
the withdrawal with that probability set to one. A limit order at $k \geq 2$ is priced $k - 1 - \priceOffset_n$ ticks inside the touch of its
side, so it moves the touch by $j$ ticks, $1 \leq j \leq k-1$, with probability
$\offsetParameter(1-\offsetParameter)^{k-1-j}$, and not at all with probability
$(1-\offsetParameter)^{k-1}$; at $k = 1$ it moves no best price.
Each entry is the probability, given the path and the earlier draws, that the two moves differ.
}
\end{frame}

\begin{frame}
  \frametitle{At one tick, two shares a touch keep the book quiet}
\begin{corol}
\label{corol.quietBook}
 At $\LOBspread = \tickSizeOfLOB$ an event moves a best price of a book iff it takes the last
 share of one of its touches.
\end{corol}

All four touches at two shares or more: the books are quiet

$\firstMove^{\cks}$, $\firstMove^{\hawOne}$: the first move of a best price of each book

Each first move: the first event emptying one of its touches
\note{
At a spread of one tick, by the table at $k = 1$: a market order or a withdrawal takes one share
from a touch, a CKS limit order at a full touch takes one share from the opposite touch, and a
geometric limit order is placed at or behind its touch. A best price moves exactly when the share
taken is the last of a touch, so with every touch at two shares or more no event moves one.
The first move of a book is the first time one of its best prices moves. Before it the spread is
the opening tick, so it is the first event that empties one of the touches of that book.
}
\end{frame}

\begin{frame}
  \frametitle{The CKS touches are walks less their running maximum}
\begin{prop}
\label{prop.cksTouchWalks}
 For $t < \firstMove^{\cks}$,
 \[
  \nthBestBidSize[1,\cks]_t = \touchWalk{b}(t) - \big(\walkMaximum(t) - \uniformDepth\big),
 \]
 likewise the ask, and $\firstMove^{\cks}$ is the first time
 $\walkMaximum - \touchWalk{b} \wedge \touchWalk{a}$ reaches $\uniformDepth$.
\end{prop}

$\touchWalk{b} = \nthBestBidSize[1]_0 + (\sideVector{b})^{\transpose}\multiCountingProc$: what
arrives at the bid touch

$\walkMaximum$: the running maximum of both walks, floored at $\uniformDepth$

$\walkMaximum - \uniformDepth$: the limit orders sent to the opposite touch
\note{
$\touchWalk{a}$ is the same with $\sideVector{a} = \unitVector_4 - \unitVector_1 - \unitVector_6$.
The proof is by induction on the events before $\firstMove^{\cks}$, before which no CKS best
price has moved, so that the spread is one tick. A limit buy raises $\touchWalk{b}$ by one, and raises $\walkMaximum$
exactly when $\touchWalk{b}$ was the maximum, that is when the CKS bid was full: the order then
trades one share at the best ask, and the bid stays while the ask falls by one, as the two
formulae say. Otherwise it joins the bid. A market sell and a withdrawal from the bid lower
$\touchWalk{b}$, leave $\walkMaximum$, and take one share from the bid. A best price first moves
at the event that takes the last share of a touch.
The cap removes the same quantity from both touches, so it couples the two sides. Under balance
the mean of each walk less its opening touch is bounded in $t$, and $\firstMove^{\cks}$ is a
function of the path.
}
\end{frame}

\begin{frame}
  \frametitle{At $\firstMove^{\cks}$ the CKS ask loses its last share}
\begin{center}
\begin{tikzpicture}[x=0.42cm, y=0.36cm, font=\scriptsize,
    bid/.style={fill=blue!40, draw=blue!60}, ask/.style={fill=orange!40, draw=orange!65}]
  \begin{scope}
    \booksizeaxis{0.4}{8.6}{4}{1}
    \draw[gray!55] (0.4,0) -- (8.6,0);
    \foreach \x in {1,...,4} \draw[bid] (\x-0.3,0) rectangle (\x+0.3,4);
    \draw[ask] (5-0.3,0) rectangle (5+0.3,1);
    \foreach \x in {6,7,8} \draw[ask] (\x-0.3,0) rectangle (\x+0.3,4);
    \node[below] at (4,-0.2) {$\bestBidPrice$};
    \node[below] at (5,-0.2) {$\bestAskPrice$};
    \node at (4.5,5.4) {before};
    \node[below, align=center] at (4.5,-1.6)
      {$\touchWalk{b} = 6$, \ $\touchWalk{a} = 3$, \ $\walkMaximum = 6$ \\[2pt]
       $\nthBestBidSize[1,\cks] = 6 - (6 - 4) = 4$ \\[2pt]
       $\nthBestAskSize[1,\cks] = 3 - (6 - 4) = 1$};
  \end{scope}
  \draw[-{Stealth}, gray!55] (9.2,2.5) -- node[above, inner sep=2pt, text=black, align=center]
    {market buy\\of one share} (14.2,2.5);
  \begin{scope}[xshift=6.3cm]
    \booksizeaxis{0.4}{8.6}{4}{1}
    \draw[gray!55] (0.4,0) -- (8.6,0);
    \foreach \x in {1,...,4} \draw[bid] (\x-0.3,0) rectangle (\x+0.3,4);
    \foreach \x in {6,7,8} \draw[ask] (\x-0.3,0) rectangle (\x+0.3,4);
    \node[below] at (4,-0.2) {$\bestBidPrice$};
    \node[below] at (6,-0.2) {$\bestAskPrice$};
    \node at (4.5,5.4) {after};
    \node[below, align=center] at (4.5,-1.6)
      {$\touchWalk{b} = 6$, \ $\touchWalk{a} = 2$, \ $\walkMaximum = 6$ \\[2pt]
       $\walkMaximum - \touchWalk{a} = 4 = \uniformDepth$};
  \end{scope}
\end{tikzpicture}
\end{center}

$\uniformDepth = 4$: the path of the next frame, at $\firstMove^{\cks}$

A market buy takes the last share of the ask

Behind the touch, $\uniformDepth$: the best ask moves one tick
\note{
The CKS book at the last event of the path of the next frame, with $\uniformDepth = 4$. Before it
the spread is one tick, the bid is full and the ask holds one share, as the proposition reads
them off the walks. A market buy lowers $\touchWalk{a}$ by one and takes that share: the drawdown
$\walkMaximum - \touchWalk{a}$ reaches $\uniformDepth$ and the ask touch empties. The best ask
moves to the next level, which holds $\uniformDepth$, as every CKS level behind a touch does: one
tick, and a spread of two. The proposition stops at this event, $\firstMove^{\cks}$.
}
\end{frame}

\begin{frame}
  \frametitle{The first move of the CKS book is a drawdown of $\uniformDepth$}
\begin{center}
\begin{tikzpicture}[x=1.05cm, y=0.3cm, font=\scriptsize]
  \path (-0.8,0) (9.3,0);
  \fill[gray!12] (0,0) -- (0,4) -- (2.6,4) -- (2.6,5) -- (4.9,5) -- (4.9,6) -- (7.4,6)
    -- (7.4,2) -- (4.9,2) -- (4.9,1) -- (2.6,1) -- (2.6,0) -- cycle;
  \draw[-{Stealth}, gray!70] (0,0) -- (8.2,0) node[right] {$t$};
  \draw[-{Stealth}, gray!70] (0,0) -- (0,7.6);
  \draw[gray!45, dashed] (0,4) -- (7.4,4);
  \node[left, gray!40!black] at (0,4) {$\uniformDepth$};
  \draw[gray!65, line width=1.6pt] plot[const plot] coordinates {
    (0,4) (2.6,5) (4.9,6) (7.4,6)};
  \draw[gray!65, dashed] plot[const plot] coordinates {
    (0,0) (2.6,1) (4.9,2) (7.4,2)};
  \draw[blue!65, thick] plot[const plot] coordinates {
    (0,3) (0.6,4) (1.4,3) (2.2,4) (2.6,5) (3.9,4) (4.4,5) (4.9,6) (5.9,5) (6.9,6) (7.4,6)};
  \draw[orange!85, thick] plot[const plot] coordinates {
    (0,2.1) (1.0,1.1) (1.8,2.1) (3.4,3.1) (5.4,4.1) (6.4,3.1) (7.4,2)};
  \fill[orange!85] (7.4,2) circle (2pt);
  \draw[gray!60!black, dashed] (7.4,0) -- (7.4,7);
  \node[below] at (7.4,0) {$\firstMove^{\cks}$};
  \node[gray!50!black, right] at (7.5,6) {$\walkMaximum$};
  \node[gray!50!black, right] at (7.5,1.6) {$\walkMaximum - \uniformDepth$};
  \node[blue!55!black, above] at (6.9,6.2) {$\touchWalk{b}$};
  \node[orange!70!black, above] at (5.9,4.1) {$\touchWalk{a}$};
\end{tikzpicture}
\end{center}

A CKS touch: the height of its walk above $\walkMaximum - \uniformDepth$
\note{
The two walks, their running maximum floored at $\uniformDepth$, and the band of width
$\uniformDepth$ below it. Each CKS touch is the height of its walk above the floor of the band, so
the walks open at the opening touches, below the cap. The bid walk raises the maximum twice: each
time the CKS bid is full and a limit buy trades at the ask, which raises the floor under the ask
walk. The ask walk meets the floor at $\firstMove^{\cks}$, where the CKS ask
empties and the best ask moves up a tick.
}
\end{frame}

\begin{frame}
  \frametitle{The CKS touches on the same path}
\begin{center}
\begin{tikzpicture}[x=1.05cm, y=0.45cm, font=\scriptsize]
  \path (-0.8,0) (9.3,0);
  \draw[-{Stealth}, gray!70] (0,0) -- (8.2,0) node[right] {$t$};
  \draw[-{Stealth}, gray!70] (0,0) -- (0,5.6);
  \draw[gray!45, dashed] (0,4) -- (7.4,4);
  \node[left, gray!40!black] at (0,4) {$\uniformDepth$};
  \foreach \t in {2.6, 4.9} {
    \draw[gray!55, dotted, thick] (\t,0) -- (\t,5);
    \node[above, gray!40!black] at (\t,5) {$\walkMaximum$ rises};}
  \draw[blue!65, thick] plot[const plot] coordinates {
    (0,3) (0.6,4) (1.4,3) (2.2,4) (3.9,3) (4.4,4) (5.9,3) (6.9,4) (7.4,4)};
  \draw[orange!85, thick] plot[const plot] coordinates {
    (0,2) (1.0,1) (1.8,2) (2.6,1) (3.4,2) (4.9,1) (5.4,2) (6.4,1) (7.4,0)};
  \fill[orange!85] (7.4,0) circle (2pt);
  \draw[gray!60!black, dashed] (7.4,0) -- (7.4,5);
  \node[below] at (7.4,0) {$\firstMove^{\cks}$};
  \node[blue!55!black, right] at (7.5,4) {$\nthBestBidSize[1,\cks]$};
  \node[orange!70!black, right] at (7.5,1) {$\nthBestAskSize[1,\cks]$};
\end{tikzpicture}
\end{center}

$\walkMaximum$ rises: a limit buy at a full bid, traded at the ask

Then the bid stays at $\uniformDepth$ and the ask loses a share

$\nthBestAskSize[1,\cks]$ reaches zero at $\firstMove^{\cks}$
\note{
The touches of the CKS book on the path of the last frame, read off the walks by the
proposition. Away from the dotted times each touch moves with its own walk: up at a limit order of
its side, down at a market order or a withdrawal that takes from it. At the dotted times the bid
is full and a limit buy arrives. Its walk rises, and so does $\walkMaximum$, but the order trades
at the best ask: the bid stays at $\uniformDepth$, and the ask loses a share at an event that is
not of its side. The ask reaches zero at $\firstMove^{\cks}$, at a market buy, and the best ask
moves up a tick.
}
\end{frame}

\begin{frame}
  \frametitle{The touch of $\hawOne$ has no cap}
\begin{prop}
\label{prop.hawkesTouch}
 For $t < \firstMove^{\hawOne}$ the best bid of $\hawOne$ is the count below, likewise the ask,
 and $\firstMove^{\hawOne}$ is the first time one of the two reaches zero.
\end{prop}
{\small
\[
 \nthBestBidSize[1,\hawOne]_t
 = \nthBestBidSize[1]_0
 + \sum_{\substack{\arrivalTimes[]_n \leq t \\ \event[n] = 3}}
 \indicator{\priceOffset_n = 0}
 - \countingProc[2](t)
 - \sum_{\substack{\arrivalTimes[]_n \leq t \\ \event[n] = 5}}
 \indicator{p_n = \bestBidPrice_{\arrivalTimes[]_n -}}
\]
}

Joins: probability $\offsetParameter$; draws the touch:
$\nthBestBidSize[1,\hawOne]/\size^{b,\hawOne}(\tickSizeOfLOB\N)$

The whole bid: $\size^{b}_0(\tickSizeOfLOB\N) = (L-1)\uniformDepth + \nthBestBidSize[1]_0$, plus a
walk
\note{
Before its first move the spread of $\hawOne$ is the opening tick, and its best bid evolves by
this count from the opening touch. A limit buy joins the touch when its offset is zero, with
probability $\offsetParameter$, and goes behind it otherwise. A market sell takes one share. A
withdrawal from the bid takes one share from the touch when it draws the touch, with probability
the touch over the whole bid, and the whole bid is the opening bid plus
$(\sideVector{b})^{\transpose}\multiCountingProc(t)$. A best price first moves at the event that
takes the last share of a touch, by the quiet book. Unlike the CKS touch, this one depends on the
draws as well as on the path.
Replace each indicator by its conditional mean and the counts by their stationary rates: that is
the drift of the next frame.
}
\end{frame}

\begin{frame}
  \frametitle{The touch of $\hawOne$ drifts, the CKS walks do not}
\begin{center}
{\small
\begin{tabular}{@{}lll@{}}
\toprule
book & drift of the bid touch & first move, of the order of \\
\midrule
$\hawOne$ & $\offsetParameter\stationaryIntensity[3] - \stationaryIntensity[2]
 - \stationaryIntensity[5]\,\nthBestBidSize[1,\hawOne]_t/\size^{b,\hawOne}_t(\tickSizeOfLOB\N)$
 & $\uniformDepth$, if negative \\
$\cks$ & $(\sideVector{b})^{\transpose}\stationaryIntensity = 0$ & $\uniformDepth^{2}$ \\
\bottomrule
\end{tabular}
}
\end{center}

Openings of the order of $\uniformDepth$; drifts of order one

Fixed $\uniformDepth$: either book may move first

$\uniformDepth \to \infty$: the two scales separate
\note{
Before their first moves the touches of the two books are walks from openings of the order of
$\uniformDepth$, and they differ in their drift. The best bid of $\hawOne$ gains a share at a
limit buy with probability $\offsetParameter$, and loses one at every market sell and at a
withdrawal that draws it. Its drift is of order one per unit of time, and below the threshold
$\offsetParameter\stationaryIntensity[3] < \stationaryIntensity[2]$ it is negative, so the touch
empties in a time of the order of $\uniformDepth$. The CKS touch walks receive every event of
their side, and under balance they have no drift: the drawdown that gives $\firstMove^{\cks}$, of
the order of $\uniformDepth$, comes of their fluctuations alone, in a time of the order of
$\uniformDepth^2$. At a fixed $\uniformDepth$ either book may move first. As $\uniformDepth$ grows
the two scales separate, and we let $\uniformDepth \to \infty$ with opening touches in proportion:
the counts follow their means over horizons proportional to $\uniformDepth$. Replace the whole
bid by the opening bid, and measure time and size in units of $\uniformDepth$: that is the fluid
equation of the next frame.
}
\end{frame}

\begin{frame}
  \frametitle{The touch of $\hawOne$ follows a fluid limit}
\begin{prop}
\label{prop.fluidTouch}
 Let $\uniformDepth \to \infty$ with $L$, $\offsetParameter$ and the specification fixed,
 $\size^{x,1}_0/\uniformDepth \to \fluidTouch^{x}(0) \in (0, 1]$, and
 $(\fluidTouch^{x})\derivative = \offsetParameter\stationaryIntensity[3] - \stationaryIntensity[2]
 - \stationaryIntensity[5]\fluidTouch^{x}/L_x$.
 On any $[0, s_1]$ where $\fluidTouch^{b}, \fluidTouch^{a} > 0$,
 $\sup_{s \leq s_1}\big\vert \size^{x,1,\hawOne}_{\uniformDepth s}/\uniformDepth
 - \fluidTouch^{x}(s)\big\vert \to 0$ in probability.
\end{prop}

Time and size in units of $\uniformDepth$; $x \in \{b, a\}$

$L_x = L - 1 + \fluidTouch^{x}(0)$: the opening side over $\uniformDepth$

$\Prob(\firstMove^{\hawOne} > \uniformDepth s_1) \to 1$
\note{
The proof is in three steps. First, a law of large numbers for the counts on the scale
$\uniformDepth$: $\multiCountingProc(\uniformDepth s)/\uniformDepth \to s\stationaryIntensity$
uniformly on $[0, s_1]$ in probability, from the identity for the count increments of class 4,
Doob's inequality for the martingale, and second moments of the state bounded in time. In
particular the whole bid over $\uniformDepth$ tends to $L_b$, the opening bid over
$\uniformDepth$.
Second, the draws: the indicators of the offsets and of the levels drawn, less their conditional
means, sum to a martingale with increments bounded by one, which is $o(\uniformDepth)$
uniformly.
Third, Gronwall's inequality for the distance between the rescaled touch and
$\fluidTouch^{b}$, which starts at the gap between the opening touch over $\uniformDepth$ and
$\fluidTouch^{b}(0)$. The ask is the same, from its own opening. Since both fluid touches stay
positive on $[0, s_1]$, so do the touches, with probability tending to one, and $\hawOne$ makes
no move before $\uniformDepth s_1$.
}
\end{frame}

\begin{frame}
  \frametitle{The fluid touch in the two regimes}
\begin{center}
\begin{tikzpicture}[x=0.68cm, y=1.35cm, font=\scriptsize]
  \draw[-{Stealth}, gray!70] (0,0) -- (12.6,0) node[right] {$s$};
  \draw[-{Stealth}, gray!70] (0,0) -- (0,2.25) node[above] {$\fluidTouch^{x}$};
  \draw[gray!45, dashed] (0,1) -- (12,1);
  \node[left, gray!40!black] at (0,1) {$\fluidTouch^{x}(0)$};
  \draw[gray!45, dashed] (0,1.6667) -- (12,1.6667);
  \node[left, gray!40!black] at (0,1.6667) {$\fluidTouch^{x}_{\infty}$};
  \draw[blue!65, thick] plot[domain=0:12, samples=60]
    (\x, {1.6667 - 0.6667*exp(-0.15*\x)});
  \draw[orange!85, thick] plot[domain=0:3.1334, samples=40]
    (\x, {-1.6667 + 2.6667*exp(-0.15*\x)});
  \fill[orange!85] (3.1334,0) circle (2pt);
  \node[below] at (3.1334,0) {$\fluidEmptying_{x}$};
  \node[blue!55!black, above left] at (12,1.6667)
       {$\offsetParameter\stationaryIntensity[3] > \stationaryIntensity[2]$};
  \node[orange!70!black, right] at (3.3,0.35)
       {$\offsetParameter\stationaryIntensity[3] < \stationaryIntensity[2]$};
\end{tikzpicture}
\end{center}

$\fluidTouch^{x}(s) = \fluidTouch^{x}_{\infty}
+ (\fluidTouch^{x}(0) - \fluidTouch^{x}_{\infty})\,e^{-\stationaryIntensity[5]s/L_x}$

$\fluidTouch^{x}_{\infty} = L_x(\offsetParameter\stationaryIntensity[3] - \stationaryIntensity[2])/\stationaryIntensity[5]$
\note{
The equation is linear: the touch relaxes at the rate $\stationaryIntensity[5]/L_x$ towards
$\fluidTouch^{x}_{\infty}$. The picture starts both curves at the same opening. Limit orders join it at the rate $\offsetParameter\stationaryIntensity[3]$, market
orders take from it at the rate $\stationaryIntensity[2]$, and withdrawals at a rate proportional to the
touch. The sign of $\offsetParameter\stationaryIntensity[3] - \stationaryIntensity[2]$ decides. Below the threshold
$\fluidTouch^{x}_{\infty}$ is negative and the fluid touch reaches zero at $\fluidEmptying_{x}$;
above it, the fluid touch settles at $\fluidTouch^{x}_{\infty} > 0$, above or below its opening
value.
}
\end{frame}

\begin{frame}
  \frametitle{Below the threshold the touch of $\hawOne$ empties first}
\begin{corol}
\label{corol.twoRegimes}
 As $\uniformDepth \to \infty$:
 \begin{enumerate}[label={\emph{\roman{*}}.}]
  \item if $\offsetParameter\stationaryIntensity[3] < \stationaryIntensity[2]$,
        $\Prob(\partingTime = \firstMove^{\hawOne}) \to 1$ and
        $\firstMove^{\hawOne}/\uniformDepth \to \fluidEmptying_{b} \wedge \fluidEmptying_{a}$
        in probability;
  \item if $\offsetParameter\stationaryIntensity[3] \geq \stationaryIntensity[2]$,
        $\partingTime/\uniformDepth \to \infty$ in probability.
 \end{enumerate}
\end{corol}

$\fluidEmptying_{x} = (L_x/\stationaryIntensity[5])
\log\big(1 + \stationaryIntensity[5]\fluidTouch^{x}(0)/L_x(\stationaryIntensity[2]
- \offsetParameter\stationaryIntensity[3])\big)$: the zero of $\fluidTouch^{x}$

Threshold: $\offsetParameter = \stationaryIntensity[2]/\stationaryIntensity[3]$

i.: the first mover and $\partingTime$; ii.: no move on the scale $\uniformDepth$
\note{
In the first regime the CKS touch walks move by $o(\uniformDepth)$ on the scale $\uniformDepth$,
and the opening touches are of order $\uniformDepth$, so their drawdown stays below
$\uniformDepth$ and the CKS book does not move. Each touch of $\hawOne$ follows its fluid touch
up to its zero, and cannot stay positive much beyond, where the continued solution is negative:
the first move is near $\uniformDepth$ times the earlier zero. There the CKS touches are still of
order $\uniformDepth$, so the CKS book moves no best price, and
$\partingTime = \firstMove^{\hawOne}$.
From the threshold up the fluid touch is positive on every $[0, s_1]$, and neither book moves on
the scale $\uniformDepth$. The proof of part ii stops there.
The regime is decided by the sign of $\offsetParameter\stationaryIntensity[3] -
\stationaryIntensity[2]$, the rate at which limit buys join the best bid of $\hawOne$ against the
rate at which market sells take from it; the withdrawals pull the touch towards
$\fluidTouch^{x}_{\infty}\uniformDepth$. In the first regime the limit decides which book moves
first, and gives the parting time to first order, as $\uniformDepth\fluidEmptying$; the CKS book
plays no part in it. At a finite $\uniformDepth$ the touch of $\hawOne$ fluctuates about the fluid
touch by an amount of the order of $\sqrt{\uniformDepth}$, so the parting time over
$\uniformDepth$ falls short of $\fluidEmptying$ on average, and by more near the threshold, where
the fluid touch meets zero at a small slope.
}
\end{frame}

\begin{frame}
  \frametitle{Beyond the corollary, $L$ orders the first moves}
\begin{center}
\begin{tabular}{@{}ll@{}}
\toprule
book & first move, of the order of \\
\midrule
$\cks$ & $\uniformDepth^{2}/(\sideVector{b})^{\transpose}\countCovarianceRate\sideVector{b}$ \\
$\hawOne$ & $(L\uniformDepth)^{2}/(\sideVector{b})^{\transpose}\countCovarianceRate\sideVector{b}$ \\
\bottomrule
\end{tabular}
\end{center}

$\cks$: a Brownian drawdown of one unit, a function of the path

$\hawOne$: a large deviation, or its side exhausted

Many levels: $\partingTime = \firstMove^{\cks}$
\note{
Above the threshold what follows is heuristic. The CKS book does not depend on
$\offsetParameter$, and it moves by the fluctuations of its touch walks. By the functional central
limit theorem for Hawkes processes of Bacry, Delattre, Hoffmann and Muzy, applied to the two walks,
whose drifts vanish under balance, they move by $\uniformDepth$ in a time of the order of
$\uniformDepth^2/(\sideVector{b})^{\transpose}\countCovarianceRate\sideVector{b}$. On that scale
$\firstMove^{\cks}$ is the drawdown of the touch walks read on a pair of correlated Brownian
motions, and it remains a function of the path.
The touch of $\hawOne$ relaxes, at the rate $\stationaryIntensity[5]/(L_x\uniformDepth)$ per
share of deviation, towards $\fluidTouch^{x}_{\infty}/L_x$ times the size of its side, which is
$L_x\uniformDepth$ plus a walk without drift, and it fluctuates about that centre by an amount of
the order of $\sqrt{\uniformDepth}$. It reaches zero by a large deviation of that fluctuation, or
when its side is nearly exhausted, which takes $L^2$ times the scale of the CKS book. Both first
moves are on the scale $\uniformDepth^2$, and which comes first is decided by $L$ and not by
$\uniformDepth$: on a book of many levels the CKS book moves first, and the parting time is a
function of the path.
Now to notebook 10, section 2, When the mid-prices part.
}
\end{frame}

\begin{frame}
  \frametitle{The discrepancy of the mid-prices}
\begin{defi}
\label{def.discrepancy}
 $\discrepancy_t := \midPriceOf{\cks}_t - \midPriceOf{\hawOne}_t$.
\end{defi}

$\discrepancy_t \in \tfrac{\tickSizeOfLOB}{2}\Z$: both mid-prices on the half-tick grid

$\discrepancy^{\hawQ}_t$: the same, with $\hawQ$ in place of $\hawOne$
\note{
We write $\discrepancy$ and not $D$, which is the rate at which the covariance of the counts
grows. The discrepancy is zero before the parting time, by the definition of
$\partingTime$, and everything that follows is about what happens after it.
}
\end{frame}

\begin{frame}
  \frametitle{Mean zero, and a variance made after $\partingTime$}
\begin{prop}
\label{prop.discrepancyMoments}
 \begin{enumerate}[label={\emph{\roman{*}}.}]
  \item $\discrepancy_t = 0$ on $\{\partingTime > t\}$;
  \item if $\nthBestBidSize[1]_0 = \nthBestAskSize[1]_0$, the law of $\discrepancy_t$ is
        symmetric about zero, and $\Expectation\discrepancy_t = 0$;
  \item then $\Variance\discrepancy_t = \Expectation[\discrepancy_t^{2};\, \partingTime \leq t]
        \geq (\tickSizeOfLOB/2)^{2}\,\Prob(\discrepancy_t \neq 0)$;
  \item $\Variance\discrepancy_t
        = \Variance\Expectation[\discrepancy_t \mid \sigmaAlgebra^{\multiCountingProc}]
        + \Expectation\Variance(\midPriceOf{\hawOne}_t \mid \sigmaAlgebra^{\multiCountingProc})$.
 \end{enumerate}
\end{prop}

$\hawQ$ in place of $\hawOne$: the same

ii.: the mirror, with the same draws $(\priceOffset_n, U_n, \orderSize_n)$, needs equal
opening touches

iv.: the flow part, then the placement part
\note{
Part i is the definition of $\partingTime$. Part ii is the mirror. The reflection
$p \mapsto 2\midPrice_0 - p$ maps the tick grid to itself and exchanges the bid and the ask; it
maps the opening book to itself when its two touches are equal, and that is where the hypothesis
is used. On $G$ both placements commute with it, the type replaced by its mirror; a
withdrawal drawn by $U_n$ reflects to the withdrawal of the other side that the same $U_n$
defines. So the books of the mirror are the reflections of those of the path, and the discrepancy
changes sign. By direction symmetry the mirror is a Hawkes process with the same specification,
so it has the law of the path, and $G$ is invariant under the exchange.
Part iii follows, a non-zero discrepancy being at least half a tick. Part iv is the law of total
variance, the CKS book being a function of the path: the flow part is what the path decides, the
placement part what the draws of $\hawOne$ add. It needs a conditioning event that is a function
of the path. With orders of one share the part of $G$ that excludes an empty side is one, the
part that excludes the floor is not, and part iv holds exactly on the first part. For $\hawQ$ the
emptying of a side depends on the sizes, and the split is a statement on the draws that keep its
sides occupied.
}
\end{frame}

\begin{frame}
  \frametitle{The mirror maps the opening book to itself}
\begin{center}
\begin{tikzpicture}[x=0.42cm, y=0.36cm, font=\scriptsize,
    bid/.style={fill=blue!40, draw=blue!60}, ask/.style={fill=orange!40, draw=orange!65},
    reflect/.style={{Stealth[length=3pt]}-{Stealth[length=3pt]}, gray!60, bend left=60},
    mirror/.style={draw=black!70, dashed, thick}]
  \begin{scope}
    \booksizeaxis{0.4}{8.6}{4}{1}
    \draw[gray!55] (0.4,0) -- (8.6,0);
    \foreach \x/\h in {1/4, 2/4, 3/4, 4/2} \draw[bid] (\x-0.3,0) rectangle (\x+0.3,\h);
    \foreach \x/\h in {5/2, 6/4, 7/4, 8/4} \draw[ask] (\x-0.3,0) rectangle (\x+0.3,\h);
    \draw[gray!60!black, dashed] (4.5,-0.4) -- (4.5,6.2);
    \foreach \x in {2,3,4} \draw[reflect] (\x,4.4) to (9-\x,4.4);
    \node[below] at (4.5,-0.4) {$\midPrice_0$};
    \node at (4.5,7.2) {equal touches};
  \end{scope}
  \begin{scope}[xshift=5.2cm]
    \booksizeaxis{0.4}{8.6}{4}{1}
    \draw[gray!55] (0.4,0) -- (8.6,0);
    \foreach \x/\h in {1/4, 2/4, 3/4, 4/3} \draw[bid] (\x-0.3,0) rectangle (\x+0.3,\h);
    \foreach \x/\h in {5/1, 6/4, 7/4, 8/4} \draw[ask] (\x-0.3,0) rectangle (\x+0.3,\h);
    \draw[mirror] (4-0.3,0) rectangle (4+0.3,1);
    \draw[mirror] (5-0.3,0) rectangle (5+0.3,3);
    \draw[gray!60!black, dashed] (4.5,-0.4) -- (4.5,6.2);
    \foreach \x in {2,3,4} \draw[reflect] (\x,4.4) to (9-\x,4.4);
    \node[below] at (4.5,-0.4) {$\midPrice_0$};
    \node at (4.5,7.2) {unequal touches};
  \end{scope}
\end{tikzpicture}
\end{center}

$p \mapsto 2\midPrice_0 - p$: the bid and the ask exchanged

Mirror path: types exchanged, the same draws $(\priceOffset_n, U_n, \orderSize_n)$

Unequal touches: the mirror opens from another book, dashed
\note{
The reflection through the opening mid maps the tick grid to itself and exchanges the bid with
the ask, level by level. On the left the touches are equal and the opening book is its own
reflection, so the path and its mirror fold from the same book: each book of the mirror is the
reflection of the book of the path, and the discrepancy changes sign. By direction symmetry the
mirror has the law of the path, and the mean of the discrepancy is zero. On the right the
reflection of the opening book, dashed, is another book, the mirror starts from it, and the
argument says nothing about the mean.
}
\end{frame}

\begin{frame}
  \frametitle{The variance of the discrepancy, by regime}
\begin{center}
{\small
\begin{tabular}{@{}lll@{}}
\toprule
regime & best prices of $\hawOne$ & rate of $\Variance(\discrepancy_t/\tickSizeOfLOB)$ \\
\midrule
second & still, short of the life of a side
 & $\pressure^{\transpose}\countCovarianceRate\pressure/(2\uniformDepth)^{2}$, all flow part \\
first & moving & a flow part and a placement part \\
\bottomrule
\end{tabular}
}
\end{center}

$2\uniformDepth\discrepancy_t/\tickSizeOfLOB$: events behind the touches, the touches, the depth
crossed

First regime: fixed placement and fixed depth, a first guess
\note{
The imbalance of any book is tied to the moves of its best prices: over a window, the imbalance
less the change at the bid touch, plus that at the ask touch, sums the moves of the best prices,
each signed and weighted by the size at its outer value. Under the assumption of the idealised
book that size is $\uniformDepth$ at every move, of one tick, and the identity is its update.
$\hawOne$ departs from it in three ways: the events it places behind its touches, which
contribute nothing to its imbalance; its touches at $t$; and queues that need not hold
$\uniformDepth$. Given the path, an event is placed behind with a probability read off the book
before it, and the first term is the sum of these, plus a martingale of known variance.
In the second regime, on horizons short against the life of a side, the best prices of $\hawOne$
seldom move, and the discrepancy is close to the displacement of the CKS mid-price, a function
of the path, at the rate in the table. In the first both mid-prices move. With a fixed placement
probability per type and a fixed depth per tick, the variance would grow linearly, its flow part
from the difference of the impacts of one event on the two mid-prices, its placement part from
the draws. Neither holds on $\hawOne$, and the split is a first guess.
Now to notebook 10, section 3, The discrepancy: mean and variance.
}
\end{frame}

\begin{frame}
  \frametitle{What random sizes change}
\[
 \underbrace{(\sideVector{b})^{\transpose}\countCovarianceRate\,\sideVector{b}}_{\hawOne}
 \quad\longrightarrow\quad
 \underbrace{(\sideVector{b})^{\transpose}\countCovarianceRate\,\sideVector{b}
 + \Variance\xi_n\,(\stationaryIntensity[3] + \stationaryIntensity[2]
 + \stationaryIntensity[5])}_{\hawQ, \ \xi_n = \orderSize_n/\meanSize}
\]

A touch moves once an order reaches its size: never quiet

Fluid touch and regimes: those of $\hawOne$, to first order

$\Expectation\discrepancy^{\hawQ}_t = 0$; the sizes add to the placement part
\note{
The display is the variance rate of the size on a side, in units of $\meanSize$: a side of
$\hawQ$ is a walk of a larger variance than a side of $\hawOne$, and it empties sooner. A market
order or a withdrawal moves a touch of $\hawQ$ whenever its size reaches the size resting there,
possibly across several levels, and the lognormal law has unbounded support, so the books are
never quiet: they may part at any event. In units the sizes have mean one, and to first order
they leave the fluid touch and the two regimes as they are. The mean of the discrepancy is still
zero, by the mirror, which leaves the sizes unchanged. Its variance gains what the count of
events does not see: the size of each order enters the placement part, and it makes random,
given the path, both the depth through which $\hawQ$ moves and its whole-book imbalance.
How fast the variance grows in each regime, how much of it the path decides, and how much the
sizes add to it, are left to measurement.
Now to notebook 10, section 4, Uncertain sizes.
}
\end{frame}
